\FigureLayoutDeclare{img/2022/national_paper_a_science/q4_fig1.png}{7.0cm}{0bp 11bp 46bp 16bp}

\chapter{2022年全国甲卷（理）}

\section{单选题}

\begin{problem}
设 \(z=-1+\sqrt{3}\i\)，则 \(\frac{z}{z\bar z-1}=\)（\quad）

\choices
    {\(-1+3\i\)}
    {\(-1-3\i\)}
    {\(-\frac13+\frac{\sqrt3}{3}\i\)}
    {\(-\frac13-\frac{\sqrt3}{3}\i\)}
\end{problem}

\begin{problem}
某社区通过公益讲座以普及社区居民的垃圾分类知识. 为了解讲座效果，随机抽取 \(10\) 位社区居民，让他们在讲座前和讲座后各回答一份垃圾分类知识问卷，这 \(10\) 位社区居民在讲座前和讲座后问卷答题的正确率如下图：

\bitmapfigure[max width=6.0cm]{img/2022/national_paper_a_science/q2_fig1.png}

则（\quad）

\choices
    {讲座前问卷答题的正确率的中位数小于 \(70\%\)}
    {讲座后问卷答题的正确率的平均数大于 \(85\%\)}
    {讲座前问卷答题的标准差小于讲座后正确率的标准差}
    {讲座后问卷答题的正确率的极差大于讲座前正确率的极差}
\end{problem}

\begin{problem}
设全集 \(U=\{-2,-1,0,1,2,3\}\)，集合 \(A=\{-1,2\}\)，\(B=\{x\mid x^2-4x+3=0\}\)，则 \(\complement_U(A\cup B)\) 为（\quad）

\choices
    {\(\{1,3\}\)}
    {\(\{0,3\}\)}
    {\(\{-2,1\}\)}
    {\(\{-2,0\}\)}
\end{problem}

\begin{problem}
如图，网格纸上绘制的是一个多面体的三视图，网格小正方形的边长为 \(1\)，则该多面体的体积为（\quad）

\bitmapfigure[width=7.0cm]{img/2022/national_paper_a_science/q4_fig1.png}

\choices
    {\(8\)}
    {\(12\)}
    {\(16\)}
    {\(20\)}
\end{problem}

\begin{problem}
函数 \(y=\left( 3^x-3^{-x} \right)\cos x\) 在区间 \(\left[ -\frac{\pi}{2},\frac{\pi}{2} \right]\) 的图象大致为（\quad）

\choices
    {\choicebitmap[width=0.15\paperwidth]{img/2022/national_paper_a_science/q5_optA.png}}
    {\choicebitmap[width=0.15\paperwidth]{img/2022/national_paper_a_science/q5_optB.png}}
    {\choicebitmap[width=0.15\paperwidth]{img/2022/national_paper_a_science/q5_optC.png}}
    {\choicebitmap[width=0.15\paperwidth]{img/2022/national_paper_a_science/q5_optD.png}}
\end{problem}

\begin{problem}
当 \(x=1\) 时，函数 \(f(x)=a\ln x+\frac b x\) 取得最大值 \(-2\)，则 \(f'(2)=\)（\quad）

\choices
    {\(-1\)}
    {\(-\frac12\)}
    {\(\frac12\)}
    {\(1\)}
\end{problem}

\begin{problem}
在长方体 \(ABCD-A_1B_1C_1D_1\) 中，已知 \(B_1D\) 与平面 \(ABCD\) 和平面 \(AA_1B_1B\) 所成角均为 \(30^\circ\)，则（\quad）

\choices
    {\(AB=2AD\)}
    {\(AB\) 与平面 \(AB_1C_1D\) 所成的角为 \(30^\circ\)}
    {\(AC=CB_1\)}
    {\(B_1D\) 与平面 \(BB_1C_1C\) 所成的角为 \(45^\circ\)}
\end{problem}

\begin{problem}
沈括的《梦溪笔谈》是中国古代科技史上的杰作，其中收录了计算圆弧长度的“会圆术”，如图，\(\myarc{AB}\) 是以 \(O\) 为圆心，\(OA\) 为半径的圆弧，\(C\) 是 \(AB\) 的中点，\(D\) 在 \(\myarc{AB}\) 上，\(CD\perp AB\)，“会圆术”给出的弧长的近似值 \(s\) 的计算公式：\(s=\overline{AB}+\frac{CD^2}{OA}\). 当 \(OA=2\)，\(\angle AOB=60^\circ\) 时，\(s=\)（\quad）

\bitmapfigure[max width=5cm]{img/2022/national_paper_a_science/q8_fig1.png}

\choices
    {\(\dfrac{11-3\sqrt3}{2}\)}
    {\(\dfrac{11-4\sqrt3}{2}\)}
    {\(\dfrac{9-3\sqrt3}{2}\)}
    {\(\dfrac{9-4\sqrt3}{2}\)}
\end{problem}

\begin{problem}
甲、乙两个圆锥的母线长相等，侧面展开图的圆心角之和为 \(2\pi\)，侧面积分别为 \(S_{\text{甲}}\) 和 \(S_{\text{乙}}\)，体积分别为 \(V_{\text{甲}}\) 和 \(V_{\text{乙}}\). 若 \(\frac{S_{\text{甲}}}{S_{\text{乙}}}=2\)，则 \(\frac{V_{\text{甲}}}{V_{\text{乙}}}=\)（\quad）

\choices
    {\(\sqrt5\)}
    {\(2\sqrt2\)}
    {\(\sqrt{10}\)}
    {\(\dfrac{5\sqrt{10}}{4}\)}
\end{problem}

\begin{problem}
椭圆 \(C:\frac{x^2}{a^2}+\frac{y^2}{b^2}=1\)（\(a>b>0\)）的左顶点为 \(A\)，点 \(P\)，\(Q\) 均在 \(C\) 上，且关于 \(y\) 轴对称. 若直线 \(AP\)，\(AQ\) 的斜率之积为 \(\frac14\)，则 \(C\) 的离心率为（\quad）

\choices
    {\( \dfrac{\sqrt3}{2}\)}
    {\( \dfrac{\sqrt2}{2}\)}
    {\( \dfrac{1}{2}\)}
    {\(\dfrac{1}{3}\)}
\end{problem}

\begin{problem}
设函数 \(f(x)=\sin\left(\omega x+\frac\pi3\right)\) 在区间 \((0,\pi)\) 上恰有三个极值点，两个零点，则实数 \(\omega\) 的取值范围是（\quad）

\choices
    {\(\left[\frac53,\frac{13}6\right)\)}
    {\(\left[\frac53,\frac{19}6\right)\)}
    {\(\left(\frac{13}6,\frac83\right]\)}
    {\(\left(\frac{13}6,\frac{19}6\right]\)}
\end{problem}

\begin{problem}
已知 \(a=\frac{31}{32}\)，\(b=\cos\frac14\)，\(c=4\sin\frac14\)，则（\quad）

\choices
    {\(c>b>a\)}
    {\(b>a>c\)}
    {\(a>b>c\)}
    {\(a>c>b\)}
\end{problem}

\section{填空题}

\begin{problem}
设向量 \(\bs a\)，\(\bs b\) 的夹角的余弦值为 \(\frac13\)，且 \(|\bs a|=1\)，\(|\bs b|=3\)，则 \((2\bs a+\bs b)\cdot\bs b=\)\fillinblank{}
\end{problem}

\begin{problem}
若双曲线 \(y^2-\frac{x^2}{m^2}=1\)（\(m>0\)）的渐近线与圆 \(x^2+y^2-4y+3=0\) 相切，则 \(m=\)\fillinblank{}
\end{problem}

\begin{problem}
从正方体的 \(8\) 个顶点中任选 \(4\) 个，则这 \(4\) 个点在同一个平面的概率为\fillinblank{}.
\end{problem}

\begin{problem}
已知 \(\triangle ABC\)，点 \(D\) 在边 \(BC\) 上，\(\angle ADB=120^\circ\)，\(AD=2\)，\(CD=2BD\). 当 \(\frac{AC}{AB}\) 取得最小值时，\(BD=\)\fillinblank{}.
\end{problem}

\section{解答题}

\begin{problem}
记 \(S_n\) 为数列 \(\{a_n\}\) 的前 \(n\) 项和. 已知 \(\frac{2S_n}{n}+n=2a_n+1\).

\begin{enumerate}
\item 证明：\(\{a_n\}\) 是等差数列；
\item 若 \(a_4\)，\(a_7\)，\(a_9\) 成等比数列，求 \(S_n\) 的最小值.
\end{enumerate}
\end{problem}

\begin{problem}
在四棱锥 \(P-ABCD\) 中，\(PD\perp\) 底面 \(ABCD\)，\(CD\parallel AB\)，\(AD=DC=CB=1\)，\(AB=2\)，\(DP=\sqrt3\).

\begin{enumerate}
\item 证明：\(BD\perp PA\)；
\item 求 \(PD\) 与平面 \(PAB\) 所成的角的正弦值.
\end{enumerate}

\bitmapfigure[max width=6.0cm]{img/2022/national_paper_a_science/q18_fig1.png}
\end{problem}

\begin{problem}
甲、乙两个学校进行体育比赛，比赛共设三个项目，每个项目胜方得 \(10\) 分，负方得 \(0\) 分，没有平局. 三个项目比赛结束后，总得分高的学校获得冠军. 已知甲学校在三个项目中获胜的概率分别为 \(0.5\)，\(0.4\)，\(0.8\)，各项目的比赛结果相互独立.

\begin{enumerate}
\item 求甲学校获得冠军的概率；
\item 用 \(X\) 表示乙学校的总得分，求 \(X\) 的分布列与期望.
\end{enumerate}
\end{problem}

\begin{problem}
设抛物线 \(C:y^2=2px\left( p>0 \right)\) 的焦点为 \(F\)，点 \(D\left( p,0 \right)\). 过 \(F\) 的直线交 \(C\) 于 \(M\)，\(N\) 两点. 当直线 \(MD\) 垂直于 \(x\) 轴时，\(|MF|=3\).

\begin{enumerate}
\item 求 \(C\) 的方程；
\item 设直线 \(MD\)，\(ND\) 与 \(C\) 的另一个交点分别为 \(A\)，\(B\)，记直线 \(MN\)，\(AB\) 的倾斜角分别为 \(\alpha\)，\(\beta\). 当 \(\alpha-\beta\) 取得最大值时，求直线 \(AB\) 的方程.
\end{enumerate}
\end{problem}

\begin{problem}
已知函数 \(f(x)=\frac{\e^x}{x}-\ln x+x-a\).

\begin{enumerate}
\item 若 \(f(x)\ge 0\)，求 \(a\) 的取值范围；
\item 证明：若 \(f(x)\) 有两个零点 \(x_1\)，\(x_2\)，则 \(x_1x_2<1\).
\end{enumerate}
\end{problem}

\section{选考题：解答题}

\begin{problem}
在直角坐标系 \(xOy\) 中，曲线 \(C_1\) 的参数方程为 \(\begin{cases}
x=\dfrac{2+t}{6}\\
y=\sqrt t
\end{cases}\)（\(t\) 为参数），曲线 \(C_2\) 的参数方程为 \(\begin{cases}
x=-\dfrac{2+s}{6}\\
y=-\sqrt s
\end{cases}\)（\(s\) 为参数）.

\begin{enumerate}
\item 写出 \(C_1\) 的普通方程；
\item 以坐标原点为极点，\(x\) 轴正半轴为极轴建立极坐标系，曲线 \(C_3\) 的极坐标方程为 \(2\cos\theta-\sin\theta=0\)，求 \(C_3\) 与 \(C_1\) 交点的直角坐标，及 \(C_3\) 与 \(C_2\) 交点的直角坐标.
\end{enumerate}
\end{problem}

\begin{problem}
已知 \(a\)，\(b\)，\(c\) 均为正数，且 \(a^2+b^2+4c^2=3\)，证明：

\begin{enumerate}
\item \(a+b+2c\le 3\)；
\item 若 \(b=2c\)，则 \(\frac1a+\frac1c\ge 3\).
\end{enumerate}
\end{problem}

